%! Conditional Probability — Unit 4, Lesson 4.5 — Student Packet Cover Sheet
\documentclass[10pt]{article}
\usepackage{apstats-article}
\usepackage{apstats-boxes}
\usepackage{ltablex}
\keepXColumns

\begin{document}

% Full-bleed navy banner
\begin{tikzpicture}[remember picture, overlay]
  \fill[navy] (current page.north west)
    rectangle ([yshift=-0.9in]current page.north east);
\end{tikzpicture}
\vspace{-0.8in}

\begin{center}
  {\color{white}\textbf{\LARGE AP Statistics}}\\[4pt]
  {\color{skymid}\large\textbf{Unit 4: Probability, Random Variables, and Probability Distributions}}\\[6pt]
  {\color{white}\textbf{\large Lesson 4.5 \quad Conditional Probability}}
\end{center}

\vspace{0.22in}
\namedateperiod
\vspace{0.18in}

\begin{learningtargetbox}
\small
\begin{enumerate}[leftmargin=1.5em, itemsep=5pt]
  \item \ldots{} calculate a conditional probability $P(A|B)$ using the formula
        $P(A|B) = P(A \cap B)/P(B)$, showing my setup clearly.
  \item \ldots{} read a two-way table to identify the joint probability and the
        conditioning event probability.
  \item \ldots{} apply the multiplication rule $P(A \cap B) = P(A) \cdot P(B|A)$
        to find a joint probability.
\end{enumerate}
\end{learningtargetbox}

\vspace{0.14in}

\begin{tocbox}
\small
\renewcommand{\arraystretch}{1.5}
\begin{tabularx}{\linewidth}{@{} c l X r @{}}
\rowcolor{sky}
\textbf{\#} & \textbf{Component} & \textbf{Description} & \textbf{Score} \\
\midrule
1 & Warm-Up        & Spiral review: joint probability and two-way tables          & \blank{1.2cm} \\
2 & Guided Notes   & Conditional probability formula; multiplication rule         & \blank{1.2cm} \\
3 & Group Activity & Sport preference $\times$ training frequency table           & \blank{1.2cm} \\
4 & Exit Ticket    & Diet and health outcome two-way table                         & \blank{1.2cm} \\
5 & Homework       & School survey: study environment $\times$ GPA range          & \blank{1.2cm} \\
\midrule
  & \multicolumn{2}{r}{\textbf{Total}} & \blank{1.2cm} \\
\end{tabularx}
\end{tocbox}

\vspace{0.14in}

\begin{remindbox}
\small
The \textbf{core idea of Lesson 4.5}: conditioning restricts the sample space to the
event that has already occurred.

\vspace{6pt}
\begin{tabularx}{\linewidth}{@{} >{\centering\arraybackslash\columncolor{goldbg}}p{0.06\linewidth}
                                    X @{}}
\textbf{\large!} &
\textbf{Formula first.}\enspace
Always write $P(A|B) = P(A \cap B) / P(B)$ before plugging in numbers.
AP exam graders reward the setup. \\[6pt]
\textbf{\large!} &
\textbf{The denominator is the conditioning event, not the grand total.}\enspace
When working from a two-way table, the denominator is the row or column total
of the given event, not the overall total. \\[6pt]
\textbf{\large!} &
\textbf{$P(A|B) \neq P(B|A)$ in general.}\enspace
The order matters --- ``the probability of a positive test given the disease''
is very different from ``the probability of the disease given a positive test.'' \\
\end{tabularx}
\end{remindbox}

\end{document}
